\documentclass[11pt]{article}
\usepackage[margin=1in]{geometry}
\usepackage{amsmath,amssymb,amsthm}
\usepackage{booktabs}
\usepackage{microtype}
\usepackage{fancyvrb}
\usepackage{hyperref}

\newtheorem{proposition}{Proposition}
\newtheorem{corollary}[proposition]{Corollary}
\theoremstyle{remark}
\newtheorem{remark}[proposition]{Remark}

\newcommand{\Norm}{\mathcal{N}}
\newcommand{\one}{\mathbf{1}}
\newcommand{\Rb}{\mathbb{R}}

\title{Single-gene Brownian motion on a lineage tree\\
\large The scBEM validation brick: model, pruning likelihood, and identifiability}
\author{scPhyTr}
\date{3 September 2026}

\begin{document}
\maketitle

\begin{abstract}
This note fixes the mathematics for the first scBEM component: one gene evolving under Brownian
motion on a fixed single-cell lineage tree, observed at the tips with independent noise. We give the
induced tip distribution, marginalise the ancestral state in closed form, derive the linear-time
pruning likelihood and prove it identical to the dense multivariate normal, and settle the
identifiability of the diffusion rate against the tip-noise variance. Two propositions and one
eigenvalue result explain precisely why the two variance components are correlated but not
degenerate, and how that depends on the shape of the tree. All numerical claims have been verified
on the two trees used as fixtures.
\end{abstract}

\section{Setup}

Let $T$ be a rooted tree with $n$ tips indexed $1,\dots,n$, internal nodes $n{+}1,\dots,2n{-}1$
(after collapsing unary chains and binarising polytomies with zero-length edges, so that $T$ is
strictly binary), root $r$, and branch length $b_u > 0$ above node $u$. Write $p(u)$ for the parent
of $u$, $d_u$ for the root-to-$u$ path length, and $\mathrm{mrca}(i,j)$ for the most recent common
ancestor of tips $i$ and $j$.

A single gene takes a latent value $x_u \in \Rb$ at every node. We observe only the tips, and only
through noise.

\section{The generative model and the induced tip distribution}

\paragraph{Brownian motion.} Conditional on its parent, each node's value is
\begin{equation}
  x_u \mid x_{p(u)} \;\sim\; \Norm\!\left(x_{p(u)},\; \sigma^2 b_u\right),
  \label{eq:bm}
\end{equation}
independently across branches, with $\sigma^2 > 0$ the diffusion rate.

\paragraph{Observation.} Tip $i$ is observed as
\begin{equation}
  y_i \mid x_i \;\sim\; \Norm\!\left(x_i,\; \tau^2\right),
  \label{eq:obs}
\end{equation}
independently across tips. The term $\tau^2$ absorbs every source of tip variation that is not
inherited: measurement and sampling noise, transcriptional bursting, and non-heritable cell state.
It is the model-internal replacement for smoothing expression over the tree as a preprocessing step.

\paragraph{Ancestral state.} The root is given a proper prior
\begin{equation}
  x_r \;\sim\; \Norm\!\left(m_0,\; v_0\right).
  \label{eq:root}
\end{equation}

Because \eqref{eq:bm}--\eqref{eq:root} are jointly Gaussian and the tips are a linear functional of
the latent field, the marginal is Gaussian and available in closed form.

\begin{proposition}[Induced tip distribution]
\label{prop:marginal}
Let $C \in \Rb^{n \times n}$ be the Brownian covariance shape,
\begin{equation}
  C_{ij} \;=\; d_{\mathrm{mrca}(i,j)}, \qquad\text{so in particular}\qquad C_{ii} = d_i .
  \label{eq:C}
\end{equation}
Then, marginally over all latent values including the root,
\begin{equation}
  y \;\sim\; \Norm\!\left(m_0 \one,\;\; \Sigma\right),
  \qquad
  \Sigma \;=\; \sigma^2 C \;+\; \tau^2 I \;+\; v_0\, \one\one^{\!\top}.
  \label{eq:Sigma}
\end{equation}
\end{proposition}

\begin{proof}
Conditional on $x_r$, the value at tip $i$ is $x_r$ plus the sum of independent increments along the
root-to-$i$ path, so $\operatorname{Cov}(x_i, x_j \mid x_r)$ is $\sigma^2$ times the length of the
shared portion of the two paths, which is the root-to-$\mathrm{mrca}(i,j)$ path, giving $\sigma^2
C_{ij}$. Adding \eqref{eq:obs} contributes $\tau^2 I$. Marginalising $x_r$ under \eqref{eq:root}
adds $v_0$ to every entry, since every tip loads on $x_r$ with coefficient one.
\end{proof}

\begin{remark}
The rank-one term $v_0 \one\one^{\!\top}$ is the part most easily dropped by accident.
Marginalising the root does not merely remove a parameter; it correlates every pair of tips.
A dense reference implementation that omits it will disagree with the pruning recursion of
Section~\ref{sec:pruning} for reasons unrelated to pruning.
\end{remark}

\section{Reparameterisation}

The pair $(\sigma^2, \tau^2)$ is a poor coordinate system for sampling, and $\sigma^2$ is not a
reportable quantity: lineage-tracing branch lengths are expressed in recorder-edit units rather than
time, so $\sigma^2$ absorbs the arbitrary scale of $C$ and is not comparable between trees. We
therefore sample the total tip variance $V$ and the phylogenetic heritability $h$,
\begin{equation}
  \sigma^2 \;=\; \frac{V h}{\bar{T}}, \qquad
  \tau^2   \;=\; V\,(1-h), \qquad
  \bar{T} \;=\; \frac{1}{n}\sum_{i=1}^n d_i ,
  \label{eq:reparam}
\end{equation}
with $V > 0$ and $h \in (0,1)$. Then $h$ is scale-free and is the quantity that transfers across
datasets, while $V$ carries the units.

\section{The pruning likelihood}
\label{sec:pruning}

Evaluating \eqref{eq:Sigma} directly costs $O(n^3)$ per likelihood evaluation. The tree makes the
precision of the latent field sparse, and Felsenstein's pruning exploits this to compute the same
number in $O(n)$ without ever forming $\Sigma$.

\paragraph{Messages.} We propagate, for each node $u$, a mean $m_u$ and a variance $V_u$ expressed
\emph{in units of $\sigma^2$}, representing the likelihood of the subtree below $u$ as a function of
$x_u$, up to a constant already accumulated:
\begin{equation}
  L_u(x_u) \;\propto\; \Norm\!\left(x_u;\; m_u,\; \sigma^2 V_u\right).
\end{equation}

\paragraph{Initialisation.} At tip $i$, \eqref{eq:obs} gives $L_i(x_i) = \Norm(x_i; y_i, \tau^2)$, so
\begin{equation}
  m_i = y_i, \qquad V_i = \tau^2/\sigma^2 .
  \label{eq:init}
\end{equation}
Heteroskedastic tip noise is accommodated at no cost by setting $V_i = \tau_i^2/\sigma^2$.

\paragraph{Propagation along a branch.} Integrating \eqref{eq:bm} over $x_u$ for a child $u$ of $w$
inflates the variance by the branch length,
\begin{equation}
  \tilde{V}_u \;=\; V_u + b_u .
  \label{eq:prop}
\end{equation}

\paragraph{Fusion at an internal node.} For an internal node $w$ with children $a$ and $b$, the
standard Gaussian product identity
\begin{equation}
  \Norm(x; \mu_1, s_1)\,\Norm(x; \mu_2, s_2)
  \;=\;
  \Norm\!\left(\mu_1; \mu_2, s_1 + s_2\right)\,
  \Norm\!\left(x; \tfrac{\mu_1 s_2 + \mu_2 s_1}{s_1+s_2}, \tfrac{s_1 s_2}{s_1+s_2}\right)
  \label{eq:gaussprod}
\end{equation}
applied with $s_1 = \sigma^2\tilde{V}_a$ and $s_2 = \sigma^2\tilde{V}_b$ gives the recursion
\begin{align}
  w_{\!\scriptscriptstyle W} &= \tilde{V}_a + \tilde{V}_b, &
  \delta_{\!\scriptscriptstyle W} &= m_a - m_b, \label{eq:contrast}\\
  V_{\!\scriptscriptstyle W} &= \frac{\tilde{V}_a \tilde{V}_b}{\tilde{V}_a + \tilde{V}_b}, &
  m_{\!\scriptscriptstyle W} &= \frac{m_a \tilde{V}_b + m_b \tilde{V}_a}{\tilde{V}_a + \tilde{V}_b},
  \label{eq:fuse}
\end{align}
and emits the factor $\Norm(\delta_{\!\scriptscriptstyle W}; 0, \sigma^2 w_{\!\scriptscriptstyle W})$,
the \emph{contrast} between the two children. The product form
$V_{\!\scriptscriptstyle W} = \tilde V_a \tilde V_b / (\tilde V_a + \tilde V_b)$ is preferred to the
harmonic form $(\tilde V_a^{-1} + \tilde V_b^{-1})^{-1}$ because it remains well defined when one of
the two vanishes, which happens whenever $\tau = 0$ meets a zero-length edge introduced by
binarising a polytomy.

\paragraph{Termination.} At the root, integrating the message against the prior \eqref{eq:root}
using \eqref{eq:gaussprod} once more emits
$\Norm\!\left(m_r - m_0;\, 0,\, \sigma^2 V_r + v_0\right)$.

\begin{proposition}[Pruning computes the dense density]
\label{prop:identity}
With $\delta$ and $w$ as in \eqref{eq:contrast} over the $n-1$ internal nodes, and $\Sigma$ as in
\eqref{eq:Sigma},
\begin{equation}
  \log \Norm\!\left(y;\, m_0\one,\, \Sigma\right)
  =
  -\frac{1}{2}\left[
    n \log(2\pi\sigma^2)
    + \sum_{W} \log w_{\!\scriptscriptstyle W} + \log v_r
    + \frac{1}{\sigma^2}\left(\sum_{W} \frac{\delta_{\!\scriptscriptstyle W}^2}{w_{\!\scriptscriptstyle W}}
      + \frac{(m_r-m_0)^2}{v_r}\right)
  \right],
  \label{eq:identity}
\end{equation}
where $v_r = V_r + v_0/\sigma^2$. In particular
$\lvert\Sigma\rvert = \sigma^{2n}\, v_r \prod_{W} w_{\!\scriptscriptstyle W}$.
\end{proposition}

\begin{proof}
Each step above is an exact Gaussian marginalisation of one latent variable, applied in post-order,
so their product telescopes to $\int p(y \mid x)\,p(x)\,dx$, which is the left-hand side. The count
of emitted factors is $(n-1) + 1 = n$, matching the dimension of $y$, and the emitted variances are
exactly the pivots of a Cholesky factorisation of $\Sigma$ in the elimination order induced by the
tree, which gives the determinant identity.
\end{proof}

Equation~\eqref{eq:identity} is the checkable statement: an implementation is correct if and only if
it reproduces the dense log-density for arbitrary $(\sigma^2, \tau^2, v_0, m_0)$.

\section{Implementing the recursion in Stan}
\label{sec:stan}

Section~\ref{sec:pruning} is an algorithm, and Stan is a declarative modelling language, so
it is worth saying precisely how one becomes the other. The answer is that the recursion is
written as an ordinary user-defined function in Stan's \texttt{functions} block, and nothing
about it is special: no external C\texttt{++}, no hand-written derivatives, no custom
distribution.

\paragraph{The tree never enters as a tree.} Stan has no recursion, no pointers and no ragged
arrays, so a post-order traversal cannot be expressed directly. Instead the tree is passed as
flat arrays, with the nodes \emph{already numbered so that the traversal is a forward loop}:
tips occupy slots $1,\dots,n$; internal nodes occupy $n{+}1,\dots,2n{-}1$; and the numbering
is chosen so that
\begin{equation}
  \text{both children of node } u \text{ have indices } < u ,
  \label{eq:ordering}
\end{equation}
which places the root last. Property \eqref{eq:ordering} is what makes
\texttt{for (k in 1:(n-1))} a valid post-order: when the loop reaches an internal node, both
of its children have already been computed. The renumbering is done once, in Python, by an
iterative depth-first walk; the tree structure reaching Stan is then two integer arrays of
child indices and one vector of branch lengths.

\paragraph{Gradients come from autodiff.} The messages $m_u$ and $V_u$ are local real-valued
variables whose values depend on the parameters, so Stan records the arithmetic of the forward
loop on its expression tape and obtains $\nabla \log p$ by a reverse sweep. Nothing about the
recursion needs to be differentiated by hand. The cost is $O(n)$ forward and $O(n)$ reverse per
leapfrog step, against $O(n^3)$ for a Cholesky factorisation of $\Sigma$.

\paragraph{The function.} Naming it with the \texttt{\_lpdf} suffix makes the sampling statement
\texttt{y \textasciitilde{} bm\_prune(...)} available as sugar for
\texttt{target += bm\_prune\_lpdf(y | ...)}. The body is a direct transcription of
\eqref{eq:init}--\eqref{eq:fuse}:

\begin{Verbatim}[fontsize=\small, xleftmargin=1em]
real bm_prune_lpdf(vector y, array[] int c1, array[] int c2, vector blen,
                   vector tau2_tip, real sigma2, real m0, real v0) {
  int n = num_elements(y);
  int N = 2 * n - 1;
  vector[N] m;   vector[N] V;              // V in units of sigma^2
  real q = 0;    real logdet = 0;
  m[1:n] = y;
  V[1:n] = tau2_tip / sigma2;
  for (k in 1:(n - 1)) {
    int a = c1[k];   int b = c2[k];
    real va = V[a] + blen[a];
    real vb = V[b] + blen[b];
    real w  = va + vb;
    q      += square(m[a] - m[b]) / w;
    logdet += log(w);
    V[n + k] = va * vb / w;
    m[n + k] = (m[a] * vb + m[b] * va) / w;
  }
  real vr = V[N] + v0 / sigma2;
  q      += square(m[N] - m0) / vr;
  logdet += log(vr);
  return -0.5 * (n * log(2 * pi() * sigma2) + logdet + q / sigma2);
}
\end{Verbatim}

Two details are deliberate. The fused variance uses the product form
$\tilde V_a \tilde V_b / (\tilde V_a + \tilde V_b)$ rather than the algebraically equivalent
harmonic form, because the latter divides by a quantity that vanishes when $\tau = 0$ meets a
zero-length edge --- precisely the situation created by binarising a polytomy. And the tip
variances arrive as a \emph{vector}, so heteroskedastic noise costs nothing; this is what a
collapsed group of $n_i$ indistinguishable cells needs, entering as $\tau^2_i = s^2/n_i$.

\paragraph{Both engines in one file.} The dense form \eqref{eq:Sigma} is written as a second
function in the same file. The \texttt{model} block uses pruning; \texttt{generated quantities}
evaluates both and returns their difference:

\begin{Verbatim}[fontsize=\small, xleftmargin=1em]
generated quantities {
  real lp_prune = bm_prune_lpdf(y | c1, c2, blen, tau2 * tip_w, sigma2, m0, v0_eff);
  real lp_dense = bm_dense_lpdf(y | C, tau2 * tip_w, sigma2, m0, v0_eff);
  real lp_diff  = lp_prune - lp_dense;
}
\end{Verbatim}

Because both are evaluated at the same draw, $\max_{\text{draws}} |\mathrm{lp\_diff}|$ is a
like-for-like test of Proposition~\ref{prop:identity} over the whole region the sampler visits,
at no extra fitting cost and with no need to match parameter values between separate runs.

\paragraph{Two optimisations available later.} Neither is applied in the first version, and
both are recorded so they are not rediscovered. First, when $\tau^2 = 0$ the variance recursion
\eqref{eq:prop}--\eqref{eq:fuse} contains no parameters at all --- $V_i = 0$ and
$\tilde V_i = b_i$ --- so it can be hoisted into \texttt{transformed data} and computed once,
leaving only the mean recursion per iteration. With homoskedastic noise it depends on the
parameters through the single scalar $\tau^2/\sigma^2$. Second, when the model is extended to
many genes the per-gene likelihoods are conditionally independent given the shared parameters,
so \texttt{reduce\_sum} partitions the gene loop across threads with near-linear speedup. That
is the property the full scBEM model relies on.

\section{Identifiability of the rate against the tip noise}

The two variance components compete: both inflate the diagonal of $\Sigma$. The following makes the
competition precise.

\begin{proposition}[Tip noise is a uniform lengthening of the terminal branches]
\label{prop:degeneracy}
Let $T_e$ denote $T$ with every terminal branch lengthened by $e > 0$, all else unchanged. Then
\begin{equation}
  C(T_e) \;=\; C(T) + e I .
\end{equation}
\end{proposition}

\begin{proof}
For $i \neq j$, $\mathrm{mrca}(i,j)$ is an internal node and the root-to-mrca path contains no
terminal branch, so $C_{ij}$ is unchanged. For $i = j$, $C_{ii} = d_i$ is the tip depth, which gains
exactly $e$.
\end{proof}

\begin{corollary}
Brownian motion with rate $\sigma^2$ and no tip noise on $T_e$ is indistinguishable from Brownian
motion with the same rate and tip noise $\tau^2 = \sigma^2 e$ on $T$. Tip noise is therefore
identified only \emph{relative to a tree taken as known and fixed}.
\end{corollary}

\begin{proposition}[Identifiability given the tree]
\label{prop:ident}
Regard $T$ as fixed. The map $(\sigma^2,\tau^2) \mapsto \sigma^2 C + \tau^2 I$ is injective unless
$C$ is a scalar multiple of $I$.
\end{proposition}

\begin{proof}
$\sigma^2 C + \tau^2 I = \sigma'^2 C + \tau'^2 I$ implies
$(\sigma^2 - \sigma'^2)\,C = (\tau'^2 - \tau^2)\, I$. If $\sigma^2 \neq \sigma'^2$ this forces
$C \propto I$; otherwise $\tau^2 = \tau'^2$.
\end{proof}

Since $C \propto I$ requires every pair of tips to share no ancestral path, injectivity holds for
every tree of interest. The practical question is conditioning, and it is answered by the smallest
eigenvalue of $C$.

\begin{proposition}[The worst-conditioned direction is a cherry contrast]
\label{prop:cherry}
Write $C = C' + \operatorname{diag}(b_1,\dots,b_n)$ where $b_i$ is the terminal branch above tip $i$
and $C'_{ij} = d_{\mathrm{mrca}(i,j)}$ for $i \neq j$, $C'_{ii} = d_{p(i)}$. Then $C'$ is the
covariance matrix of the parent values $(x_{p(1)},\dots,x_{p(n)})$ under \eqref{eq:bm} with
$\sigma^2 = 1$, hence positive semidefinite; and it is singular whenever $T$ contains a
\emph{cherry}, meaning two tips that share a parent. Consequently, if all terminal branches have
the common length $b$,
\begin{equation}
  \lambda_{\min}(C) = b,
\end{equation}
attained on the contrast $e_i - e_j$ between the two tips of a cherry. Along that direction
$\sigma^2 C + \tau^2 I$ acts as multiplication by $\sigma^2 b + \tau^2$, so the data constrain only
that combination.
\end{proposition}

\begin{proof}
$C'$ is a Gram matrix of the parent values by the same argument as
Proposition~\ref{prop:marginal}, and if tips $i,j$ share a parent then rows $i$ and $j$ of $C'$
coincide, so $(e_i - e_j) \in \ker C'$. With $b_i \equiv b$ we have $C = C' + bI$, whence
$\lambda_{\min}(C) = \lambda_{\min}(C') + b = b$ with the same eigenvector.
\end{proof}

\begin{remark}[How bad it is, in numbers]
The dimension of the near-degenerate subspace grows with the number of cherries, so the effect is not
a single awkward direction but a substantial part of the spectrum. The Fisher information for
$(\sigma^2,\tau^2)$ has entries
$\mathcal{I}_{11} = \tfrac12\operatorname{tr}[(\Sigma^{-1}C)^2]$,
$\mathcal{I}_{22} = \tfrac12\operatorname{tr}[\Sigma^{-2}]$,
$\mathcal{I}_{12} = \tfrac12\operatorname{tr}[\Sigma^{-1}C\Sigma^{-1}]$,
and the implied correlation is reported in Table~\ref{tab:fixtures}. It sits near $-0.74$ on both
fixtures: strongly correlated, comfortably short of degenerate.
\end{remark}

\begin{table}[t]
\centering
\begin{tabular}{lrr}
\toprule
 & \emph{C.\ elegans} & SCOUT toy (casNJ) \\
\midrule
tips $n$                                   & 391      & 256 \\
nodes after binarisation                   & 781      & 511 \\
terminal branch lengths                    & 1--4     & all 1 \\
tip depths                                 & 5--11    & 6--9 \\
cherries                                   & 174      & 124 \\
$\lambda_{\min}(C)$                        & 1.000    & 1.000 \\
$\lambda_{\max}(C)$                        & 484.7    & 257.7 \\
Fisher $\operatorname{corr}(\sigma^2,\tau^2)$ at $\sigma^2=\tau^2=1$ & $-0.741$ & $-0.737$ \\
\bottomrule
\end{tabular}
\caption{The two fixtures, measured after collapsing unary chains and binarising polytomies.
$\lambda_{\min}(C)$ equals the shortest terminal branch in both cases, as
Proposition~\ref{prop:cherry} predicts. Note that the two trees are \emph{not} appreciably different
in conditioning despite their different depth spreads.}
\label{tab:fixtures}
\end{table}

\section{An exactly conjugate test configuration}

Validating a sampler against another sampler is weak. The following configuration admits a closed
form and therefore permits a check against a number.

\begin{proposition}[Exact posterior]
Set $\tau^2 = 0$ and scale the root prior with the rate, $v_0 = \sigma^2 k$ for fixed $k > 0$. Then
with $M = C + k\one\one^{\!\top}$,
\begin{equation}
  y \mid \sigma^2 \sim \Norm\!\left(m_0\one,\; \sigma^2 M\right),
  \qquad
  Q = (y - m_0\one)^{\!\top} M^{-1} (y - m_0\one),
\end{equation}
and under $\sigma^2 \sim \mathrm{InvGamma}(a_0,b_0)$ the posterior is exactly
\begin{equation}
  \sigma^2 \mid y \;\sim\; \mathrm{InvGamma}\!\left(a_0 + \tfrac{n}{2},\; b_0 + \tfrac{Q}{2}\right),
  \label{eq:conjugate}
\end{equation}
with generalised least squares estimator $\hat{\sigma}^2 = Q/n$.
\end{proposition}

The production model uses an unscaled $v_0$, for which no conjugacy holds; \eqref{eq:conjugate} is a
test configuration selected by a flag, not the model.

\section{Validation protocol}

\begin{enumerate}
  \item \textbf{Simulator agreement.} Two independent simulators --- Cholesky of \eqref{eq:Sigma},
        and a recursive walk down the tree implementing \eqref{eq:bm}--\eqref{eq:obs} --- must agree
        in distribution before either is used.
  \item \textbf{Engine agreement.} The dense form \eqref{eq:Sigma} and the recursion
        \eqref{eq:identity} must agree to $10^{-8}$ over a grid in
        $(\sigma^2,\tau^2,v_0)$. A \textsc{numpy} prototype already agrees to $\sim 10^{-12}$ on both
        fixtures across 27 settings.
  \item \textbf{Exact arithmetic.} In the conjugate configuration, posterior quantiles must match
        \eqref{eq:conjugate} within Monte Carlo error.
  \item \textbf{Recovery and calibration.} The truth's posterior quantile must be uniform across
        repeated simulations, and simulation-based calibration rank histograms must pass a
        uniformity test at the $1\%$ level over 200 replicates.
  \item \textbf{Scaling.} Wall-clock for the dense and pruning engines as $n$ grows, reported rather
        than thresholded.
\end{enumerate}

Propositions~\ref{prop:degeneracy}, \ref{prop:ident} and \ref{prop:cherry}, and the engine agreement
in item 2, have been verified numerically on both fixtures; the remaining items require Stan and have
not been run.

\end{document}
